\section{Doctrine, discipline, prudence, and fact}

Doctrine teaches what is true. Discipline orders ecclesial action. The same
document may contain both, and a discipline may embody doctrine without being
irreformable in its concrete form. Fasting rules, canonical procedures,
liturgical calendars, and governance arrangements can bind while remaining
changeable. Disobedience is not excused merely by changeability; irreformability
is not proved merely by obligation.

Prudential judgments join stable principles to mutable facts: a diplomatic
strategy, an assessment of danger, an administrative restriction, or the
credibility of an alleged apparition. They deserve reception proportionate to
office and competence, but they do not receive supernatural guarantees
foreign to their object. Historical assertions used in magisterial documents
retain their historical character unless necessary to the defined object.

The ``hierarchy of truths'' in \emph{Unitatis redintegratio} 11 is another
distinction. It orders Catholic doctrines by their relation to the foundation
of Christian faith; it does not rank some truths as false, dispensable, or
unauthoritative. A truth may be remote from the center yet definitively taught.
Centrality, certainty, and pedagogical priority are related questions, not one
scale.

